% Auto-generated complete chapter from 19B_github_readme_twinsight_final.md
\chapter{19B\_github\_readme\_twinsight\_final}
\label{complete:root_022}

\begin{sourcemeta}
\textbf{Fuente:} \href{file:///E:/Laboral/19B_github_readme_twinsight_final.md}{19B\_github\_readme\_twinsight\_final.md}\\
\textbf{Seccion:} root\\
\textbf{Estado:} canonical\\
\textbf{Rol:} Modulo final complementario\\
\textbf{Lineas originales:} 787\\
\textbf{Ultima modificacion:} 2026-06-11 19:52\\
\textbf{Deduplicacion conservadora:} 0 bloques repetidos omitidos; 0 caracteres repetidos no reimpresos.
\end{sourcemeta}

\section*{Resumen de orientacion}
This file provides a final professional README structure and copy for the GitHub repository associated with:

\section*{Contenido completo depurado}
\addcontentsline{toc}{section}{Contenido completo depurado}




\subsection{0. Module status}




\begin{mdcode}
Bloque de codigo: text
Module: 19B_github_readme_twinsight_final.md
Date generated: 2026-06-12
Status: Complete v1
Purpose: Final GitHub README content for TwinSight X500 / WebGL-Thesis-Proposal
Inputs: 08B, 19, 20B, 21B, 29, 29B, 29C
Use: GitHub README, portfolio technical proof, interview review, recruiter validation
\end{mdcode}




\medskip\hrule\medskip




\section{1. Purpose}




This file provides a final professional README structure and copy for the GitHub repository associated with:




\begin{mdcode}
Bloque de codigo: text
TwinSight X500
Unity WebGL Technical Visualization for Drone Assembly Inspection
\end{mdcode}




The README must support Alexander’s positioning as:




\begin{mdcode}
Bloque de codigo: text
Real-Time 3D Developer / Unity Technical Artist
\end{mdcode}




with evidence in:




\begin{mdcode}
Bloque de codigo: text
Unity · C# · WebGL · URP · UI Toolkit · Blender · CAD-to-Realtime · Technical Visualization
\end{mdcode}




The README should be concise enough for recruiters and structured enough for technical reviewers.




\medskip\hrule\medskip




\section{2. README goals}




\subsection{2.1 Primary goals}




\begin{mdcode}
Bloque de codigo: text
1. Explain the project in under 10 seconds.
2. Show visual proof immediately.
3. Provide demo/video/portfolio links.
4. Clarify technical stack.
5. Explain the problem and solution.
6. Show key features.
7. Document CAD-to-realtime optimization.
8. Present evaluation metrics.
9. State limitations honestly.
10. Clarify AI-assisted workflow if relevant.
\end{mdcode}




\subsection{2.2 What the README should not do}




Avoid turning the README into:




\begin{mdcode}
Bloque de codigo: text
a full thesis document
a generic academic abstract
a long personal story
a marketing-only page
a code dump without visual proof
\end{mdcode}




\medskip\hrule\medskip




\section{3. Final README skeleton}




Use this structure in the actual repository:




\begin{mdcode}
Bloque de codigo: markdown
# TwinSight X500

Unity WebGL technical visualization prototype for drone assembly inspection.

[Demo thumbnail / video link]

## Overview
## Demo
## Key Features
## Screenshots
## Technical Stack
## Problem
## Solution
## Technical Pipeline
## CAD-to-Realtime Optimization
## Unity WebGL Architecture
## Evaluation
## My Role
## AI-Assisted Workflow
## Limitations
## Roadmap
## Repository Structure
## Setup / Running Locally
## Links
## License / Credits
\end{mdcode}




\medskip\hrule\medskip




\section{4. Final README copy}




The following block can be copied into the repository and adapted with final links.




\begin{mdcode}
Bloque de codigo: markdown
# TwinSight X500

**Unity WebGL technical visualization prototype for drone assembly inspection.**

TwinSight X500 is an interactive browser-based 3D viewer for inspecting the assembly of a Holybro X500 V2 drone. It transforms CAD-derived drone assets into an optimized real-time Unity WebGL experience with component selection, exploded view, cross-section tools, visual modes and technical information panels.

> From static 2D documentation to interactive spatial understanding.

---

## Demo

[![TwinSight X500 Demo](./media/twinsight-demo-thumbnail.jpg)](https://your-demo-video-link.com)

- **Live WebGL demo:** [add link]
- **Portfolio case study:** [add link]
- **Technical breakdown:** [add link]
- **ArtStation breakdown:** [add link]

---

## Highlights

- Unity WebGL technical visualization prototype.
- Browser-based drone assembly inspection.
- CAD-derived geometry optimized for real-time interaction.
- Component-level selection and highlighting.
- Technical part information panels.
- Exploded assembly view.
- Cross-section / clipping inspection.
- Multiple visual modes: realistic, X-ray, ghosted, blueprint, wireframe, solid and thermal-style.
- Formative usability and workload evaluation using SUS, NASA-TLX Raw and Think-Aloud.

---

## Screenshots

| Viewer | Exploded View | Cross-Section |
|---|---|---|
| ![TwinSight viewer](./media/01_viewer.jpg) | ![Exploded view](./media/02_exploded_view.jpg) | ![Cross-section](./media/03_cross_section.jpg) |

| Visual Modes | Technical UI | CAD Optimization |
|---|---|---|
| ![Visual modes](./media/04_visual_modes_grid.jpg) | ![Technical UI](./media/05_technical_ui.jpg) | ![CAD optimization](./media/06_cad_optimization.jpg) |

---

## Technical Stack

| Area | Tools / Technologies |
|---|---|
| Engine | Unity |
| Programming | C# |
| Deployment | Unity WebGL |
| Rendering | URP, Shader Graph |
| UI | UI Toolkit |
| 3D pipeline | Blender |
| Asset pipeline | CAD-derived geometry, mesh cleanup, optimization |
| Evaluation | SUS, NASA-TLX Raw, Think-Aloud |

---

## Problem

Drone assembly documentation is often presented through static 2D diagrams, part lists and written instructions. This forces users to mentally reconstruct spatial relationships between components, orientations, fasteners and hidden structures.

For complex assemblies such as a Holybro X500 V2 drone, this increases cognitive effort and makes it harder to understand how parts relate to each other in 3D space.

---

## Solution

TwinSight X500 converts CAD-derived drone geometry into an optimized Unity WebGL viewer. Instead of reading static documentation, users can navigate the assembly, select components, inspect technical information, activate exploded view, cut through the model and switch between visual modes.

The goal is not only to show the drone, but to make the relationships between its parts easier to understand.

---

## Key Features

### Component Selection

Users can select drone components directly in the 3D viewer. The selected part is highlighted and linked to contextual technical information.

### Exploded View

The exploded view separates assemblies while preserving spatial logic, helping users understand component hierarchy and placement.

### Cross-Section / Clipping

The clipping tool cuts through the model to reveal hidden structures and internal relationships.

### Visual Modes

TwinSight includes several inspection modes:

- Realistic
- X-ray
- Ghosted
- Blueprint
- Wireframe
- Solid
- Thermal-style

> The thermal-style mode is a qualitative visual mode. It is not a physical thermal simulation, FEA result or sensor-driven thermal analysis.

### Technical UI

The interface provides technical part information while preserving the 3D inspection context.

---

## Technical Pipeline

\end{mdcode}


CAD-derived source

        ↓

Tessellation / conversion route exploration

        ↓

Blender cleanup and optimization

        ↓

Selectable component organization

        ↓

Unity import and scene setup

        ↓

Runtime interaction systems

        ↓

UI Toolkit interface

        ↓

Unity WebGL build

        ↓

Usability and workload evaluation


\begin{mdcode}
Bloque de codigo

The pipeline focused on preserving the visual information needed for technical inspection while reducing complexity enough for browser-based real-time interaction.

---

## CAD-to-Realtime Optimization

A central technical challenge was converting complex CAD-derived geometry into a model suitable for Unity WebGL.

| Metric | Value |
|---|---:|
| Optimized geometry | 95,617 triangles |
| Source CAD routes | 6.5M+ triangles |
| Canonical research parts | 28 |
| Scene nodes / anchors | 30 |
| Technical elements | 257 |

The final optimized model was prepared for browser-based real-time inspection while preserving the component structure needed for selection, exploded view and technical visualization.

---

## Unity WebGL Architecture

The Unity application coordinates runtime selection, visualization states, UI panels and inspection tools inside a WebGL build.

Main runtime systems:

- Model hierarchy and selectable parts.
- Selection and highlight system.
- Camera and navigation logic.
- Exploded-view state controller.
- Clipping / cross-section controller.
- Visual-mode switching system.
- Technical metadata / part information.
- UI Toolkit panels.
- WebGL deployment configuration.

---

## Evaluation

TwinSight X500 was evaluated as an academic/portfolio prototype using usability, workload and qualitative observation methods.

| Metric | Value |
|---|---:|
| Participants | 12 |
| Task-condition records | 96 |
| SUS average | 91.88 |
| NASA-TLX Raw, 3D viewer | 8.69 |
| NASA-TLX Raw, 2D support | 19.89 |

The formative evaluation suggested high perceived usability and lower perceived workload for the 3D viewer condition compared with 2D support in the tested tasks.

---

## My Role

I worked across the complete prototype pipeline:

- Prepared and optimized CAD-derived drone geometry for real-time WebGL use.
- Organized selectable components and scene hierarchy.
- Implemented runtime interaction systems in Unity/C#.
- Built component selection, highlighting and information panels.
- Developed exploded view and clipping/cross-section interaction.
- Implemented multiple visual inspection modes.
- Designed technical UI and mobile-first interaction logic.
- Prepared evaluation workflow using SUS, NASA-TLX Raw and Think-Aloud.
- Documented constraints, findings and future improvements.

The source CAD model was used as the technical basis for the drone geometry. My work focused on transforming, optimizing, organizing and implementing it as an interactive real-time WebGL visualization prototype.

---

## AI-Assisted Workflow

AI assistance was used as an implementation and iteration support tool during parts of the development process. Final architecture decisions, debugging, Unity integration, asset preparation, evaluation design, documentation and project ownership remained my responsibility.

---

## Limitations

TwinSight X500 is an academic/portfolio prototype, not a deployed industrial product.

Current limitations:

- Not a production digital twin.
- No live IoT or sensor data.
- No predictive maintenance.
- No WebAR placement.
- No real thermal simulation.
- Formative academic evaluation.
- Prototype scope and WebGL constraints.

---

## Roadmap

Possible future improvements:

- Guided assembly sequences.
- Deeper metadata for each component.
- Search and part filtering.
- WebGPU or performance-focused rendering experiments.
- Optional XR inspection mode.
- Sensor/state data as a digital-twin extension.
- Expansion to other drones or technical assemblies.
- Improved mobile interaction and loading performance.

---

## Repository Structure

\end{mdcode}


.

├── Assets/

│   ├── Scripts/

│   ├── Materials/

│   ├── Models/

│   ├── UI/

│   └── Scenes/

├── ProjectSettings/

├── Packages/

├── media/

│   ├── twinsight-demo-thumbnail.jpg

│   ├── 01\_viewer.jpg

│   ├── 02\_exploded\_view.jpg

│   ├── 03\_cross\_section.jpg

│   ├── 04\_visual\_modes\_grid.jpg

│   ├── 05\_technical\_ui.jpg

│   └── 06\_cad\_optimization.jpg

└── README.md


\begin{mdcode}
Bloque de codigo

Adapt this structure to the actual repository before publishing.

---

## Setup / Running Locally

> Add final setup instructions after confirming the repository structure.

Example:

\end{mdcode}


git clone \url{https://github.com/your-user/your-repo.git}


\begin{mdcode}
Bloque de codigo

Open the project in Unity:

\end{mdcode}


Unity version: [add version]

Target platform: WebGL

Render pipeline: URP


\begin{mdcode}
Bloque de codigo

Build instructions:

\end{mdcode}


File → Build Settings → WebGL → Build


\begin{mdcode}
Bloque de codigo

If the repository does not include full source assets or large media files, state that clearly.

---

## Links

- Portfolio: [add link]
- Live WebGL demo: [add link]
- Demo video: [add link]
- ArtStation breakdown: [add link]
- LinkedIn: [add link]

---

## License and Credits

Source drone model:

\end{mdcode}


Holybro X500 V2 source/reference model used as the basis for technical visualization.


\begin{mdcode}
Bloque de codigo

Add exact attribution and license terms before publishing.

Project:

\end{mdcode}


TwinSight X500 developed by Alexander Woodcock Salomón as an academic/portfolio prototype.


\begin{mdcode}
Bloque de codigo

\end{mdcode}




\medskip\hrule\medskip




\section{5. Required media folder}




Create:




\begin{mdcode}
Bloque de codigo: text
/media
\end{mdcode}




Recommended files:




\begin{mdcode}
Bloque de codigo: text
twinsight-demo-thumbnail.jpg
01_viewer.jpg
02_exploded_view.jpg
03_cross_section.jpg
04_visual_modes_grid.jpg
05_technical_ui.jpg
06_cad_optimization.jpg
07_pipeline_diagram.jpg
08_metrics_card.jpg
09_architecture_diagram.jpg
10_browser_webgl_view.jpg
\end{mdcode}




\subsection{5.1 Minimum acceptable README media}




\begin{mdcode}
Bloque de codigo: text
[ ] demo thumbnail
[ ] final viewer screenshot
[ ] exploded view screenshot
[ ] cross-section screenshot
[ ] visual modes grid
[ ] optimization before/after
[ ] metrics card
\end{mdcode}




\subsection{5.2 Strong README media}




\begin{mdcode}
Bloque de codigo: text
[ ] 90-second demo video thumbnail
[ ] 30-second GIF preview
[ ] WebGL browser screenshot
[ ] Unity editor screenshot
[ ] Blender cleanup screenshot
[ ] visual modes grid
[ ] SUS/NASA-TLX card
[ ] CAD-to-realtime pipeline diagram
\end{mdcode}




\medskip\hrule\medskip




\section{6. Badge recommendations}




Use badges sparingly.




Recommended badges:




\begin{mdcode}
Bloque de codigo: markdown
![Unity](https://img.shields.io/badge/Unity-000000?logo=unity&logoColor=white)
![C#](https://img.shields.io/badge/C%23-239120?logo=csharp&logoColor=white)
![WebGL](https://img.shields.io/badge/WebGL-990000?logo=webgl&logoColor=white)
![URP](https://img.shields.io/badge/URP-Unity-blue)
![Status](https://img.shields.io/badge/status-portfolio%20prototype-orange)
\end{mdcode}




Avoid excessive badges.




Do not use badges that imply production deployment, enterprise use or certification.




\medskip\hrule\medskip




\section{7. GitHub repository description}




Use this in the repository short description:




\begin{mdcode}
Bloque de codigo: text
Unity WebGL technical visualization prototype for drone assembly inspection, featuring CAD-to-realtime optimization, component selection, exploded view, clipping tools and visual modes.
\end{mdcode}




If too long:




\begin{mdcode}
Bloque de codigo: text
Unity WebGL technical visualization prototype for drone assembly inspection.
\end{mdcode}




\medskip\hrule\medskip




\section{8. GitHub topics}




Recommended topics:




\begin{mdcode}
Bloque de codigo: text
unity
unity-webgl
csharp
webgl
technical-visualization
real-time-3d
technical-art
cad-optimization
interactive-3d
blender
drone
ui-toolkit
urp
\end{mdcode}




Optional:




\begin{mdcode}
Bloque de codigo: text
digital-twin
simulation
xr
usability-testing
nasa-tlx
sus
\end{mdcode}




Use \texttt{digital-twin} only if README clearly states that the project is digital-twin-adjacent, not a full production digital twin.




\medskip\hrule\medskip




\section{9. README visual order}




Recommended order:




\begin{mdcode}
Bloque de codigo: text
1. Demo thumbnail
2. Highlights
3. Screenshot grid
4. Technical stack
5. Problem
6. Solution
7. Features
8. Pipeline
9. Optimization
10. Architecture
11. Evaluation
12. Role
13. Limitations
14. Setup
15. Links
\end{mdcode}




Reason:




Recruiters and technical reviewers should see proof before reading details.




\medskip\hrule\medskip




\section{10. Claim-safety rules}




\subsection{10.1 Safe wording}




Use:




\begin{mdcode}
Bloque de codigo: text
technical visualization prototype
browser-based 3D viewer
CAD-derived geometry
optimized for WebGL
formative evaluation
digital-twin-adjacent
qualitative thermal-style visual mode
AI-assisted implementation support
\end{mdcode}




\subsection{10.2 Avoid}




\begin{mdcode}
Bloque de codigo: text
industrial-grade digital twin
production-ready drone twin
real thermal simulation
predictive maintenance system
IoT-integrated platform
enterprise deployment
senior-level Unity system
fully automated CAD pipeline
\end{mdcode}




Unless those claims become true and are supported.




\medskip\hrule\medskip




\section{11. README variants}




\subsection{11.1 Recruiter-friendly shorter README}




Use if the repository should be brief:




\begin{mdcode}
Bloque de codigo: markdown
# TwinSight X500

Unity WebGL technical visualization prototype for drone assembly inspection.

TwinSight X500 transforms CAD-derived drone geometry into an optimized browser-based 3D viewer with component selection, exploded view, cross-section tools, visual modes and technical information panels.

## Demo

[Demo video]
[Live WebGL build]
[Portfolio case study]

## Highlights

- Unity WebGL
- C#
- CAD-to-realtime optimization
- Component selection
- Exploded view
- Clipping tools
- Technical UI
- SUS / NASA-TLX evaluation

## Metrics

| Metric | Value |
|---|---:|
| Optimized triangles | 95,617 |
| Source CAD routes | 6.5M+ |
| SUS average | 91.88 |
| Participants | 12 |

## Stack

Unity, C#, WebGL, URP, UI Toolkit, Blender.

## My Role

I developed the Unity WebGL integration, runtime interaction systems, technical UI, visual modes, CAD-to-realtime optimization workflow, documentation and evaluation pipeline.

## Links

- Portfolio:
- Demo:
- Video:
\end{mdcode}




\subsection{11.2 Technical longer README}




Use the full version from section 4.




\medskip\hrule\medskip




\section{12. README update checklist}




\begin{mdcode}
Bloque de codigo: text
[ ] Repository description updated
[ ] Topics added
[ ] README title updated
[ ] Demo thumbnail added
[ ] Demo video linked
[ ] Live WebGL link added
[ ] Portfolio case study linked
[ ] Screenshot grid added
[ ] Tech stack table added
[ ] Feature list added
[ ] Pipeline diagram added
[ ] Optimization metrics added
[ ] Evaluation metrics added
[ ] Role/contribution clarified
[ ] AI-assisted workflow note added
[ ] Limitations added
[ ] Setup instructions verified
[ ] Credits/license clarified
[ ] Broken links checked
\end{mdcode}




\medskip\hrule\medskip




\section{13. Missing information before final publication}




Before publishing the README, fill:




\begin{mdcode}
Bloque de codigo: text
[ ] final repository URL
[ ] live WebGL build URL
[ ] YouTube/Vimeo demo URL
[ ] portfolio case-study URL
[ ] ArtStation breakdown URL
[ ] final Unity version
[ ] exact source model attribution/license
[ ] whether Marmoset/RizomUV were actually used
[ ] whether public build includes all assets
[ ] final setup instructions
\end{mdcode}




\medskip\hrule\medskip




\section{14. Integration with portfolio}




The README should link outward:




\begin{mdcode}
Bloque de codigo: text
GitHub → Portfolio case study
GitHub → Demo video
GitHub → Live WebGL build
GitHub → ArtStation breakdown
\end{mdcode}




The portfolio should link inward:




\begin{mdcode}
Bloque de codigo: text
Portfolio → GitHub
Portfolio → Live demo
Portfolio → Video
Portfolio → ArtStation
\end{mdcode}




This creates a proof loop.




\medskip\hrule\medskip




\section{15. Integration with LinkedIn}




Use README as technical proof, but do not send recruiters only to GitHub.




Preferred recruiter flow:




\begin{mdcode}
Bloque de codigo: text
LinkedIn → Portfolio case study → Demo video → GitHub
\end{mdcode}




Developer flow:




\begin{mdcode}
Bloque de codigo: text
LinkedIn → GitHub → Demo video → Technical breakdown
\end{mdcode}




\medskip\hrule\medskip




\section{16. Next module}




Proceed with:




\begin{mdcode}
Bloque de codigo: text
30_live_job_postings_market_snapshot.md
\end{mdcode}




\subsection{Tool recommendation}




\begin{center}
\scriptsize
\begin{longtable}{>{\raggedright\arraybackslash}p{0.455\linewidth} >{\raggedright\arraybackslash}p{0.455\linewidth}}
\toprule
Tool & Use \\
\midrule
\endfirsthead
\toprule
Tool & Use \\
\midrule
\endhead
Agent mode & No, unless adding more jobs \\
Deep Research & No \\
\bottomrule
\end{longtable}
\end{center}




Reason:




B1 and B5 already exist. The next useful step is not more browsing; it is to integrate B1, B5 and the B2 partial capture into a cleaned market snapshot with gaps, role signals and portfolio implications.



